% 364_Ikosaeder_D4_Vergleich_De_ch.tex
% Johann Pascher, 14. Sept. 2026
% Der Ikosaeder-Träger und das D4-Gitter: Galois-Struktur, Reduktion mod 3 und die Verdopplung 3+3'

\providecommand{\BM}{\textbf{[B]}}
\providecommand{\KM}{\textbf{[K]}}
\providecommand{\SM}{\textbf{[S]}}
\providecommand{\SetzM}{\textbf{[SETZUNG]}}
\providecommand{\QM}{\textbf{[Q]}}
\providecommand{\EM}{\textbf{[E]}}
\providecommand{\XM}{\textbf{[X]}}

\chapter*{Der Ikosaeder-Träger und das D\textsubscript{4}-Gitter}
\addcontentsline{toc}{chapter}{Der Ikosaeder-Träger und das D₄-Gitter}

\begin{abstract}
Ein extern vorgeschlagenes Trägermodell für emergente Physik verwendet einen
Zwölf-Port-Ikosaeder mit Symmetriegruppe $A_5$ und einer lokalen
Mittelungs-Reparaturregel auf den dreißig Kanten. Dieses Dokument vergleicht
diesen Träger exakt mit der FFGFT-Geometrie $T^4/\mathbb{Z}_3$ auf dem
$D_4$-Gitter. Vier Prüfskripte belegen: (i)~das Spektrum des Trägers liegt
in $\mathbb{Q}(\sqrt5)$, die Galois-Konjugation $\sqrt5\mapsto-\sqrt5$
vertauscht die beiden dreidimensionalen $A_5$-Blöcke $3\leftrightarrow3'$
als Spektralprojektoren, so dass über $\mathbb{Q}$ nur der 6D-Block
existiert -- die Verdopplung $6=3\oplus3'$ aus Dok.~285; (ii)~die
Reduktion $\mathbb{Z}[\varphi]/(3)=GF(9)$ überführt diese Galois-Konjugation
in den Frobenius $x\mapsto x^3$ des FFGFT-Grundkörpers; (iii)~$A_5$ wirkt
treu und absolut irreduzibel auf $GF(9)^3$, aber nicht auf $GF(9)^2$,
also nicht in $G(3)$, wohl aber in $G(6)$ -- dieselbe Schwelle wie die
Leptonen-Ordnung~26 in Dok.~341; (iv)~$A_5$ ist mit dem $D_4$-Gitter
unverträglich ($5\nmid1152$, Dok.~314), der gemeinsame Kern ist die
Tetraedergruppe $A_4$, die einfach transitiv auf den zwölf Ports wirkt,
und die Port-$\mathbb{Z}_3$ ist nachweislich nie eine der sechzehn
Orbifold-Erzeugenden von $T^4/\mathbb{Z}_3$. Die Eichzuordnung des
externen Modells ($u(1)\oplus su(2)\oplus su(3)$ auf $1\oplus3\oplus(3'\oplus5)$)
ist darstellungstheoretisch konsistent; ihre Lie-Struktur ist dort ein
Axiom, nicht eine Folge des Ikosaeders.
\end{abstract}

% ============================================================
\section*{Statusmarkierungen}
% ============================================================
\begin{center}
\begin{tabular}{@{}lp{10cm}@{}}
\textbf{[B]} & Algebraisch bewiesen (Prüfskript, exakte Arithmetik). \\[3pt]
\textbf{[K]} & Numerisch bestätigt. \\[3pt]
\textbf{[S]} & Strukturell plausibel, physikalische Identifikation offen. \\[3pt]
\textbf{[E]} & Etablierte externe Mathematik. \\[3pt]
\textbf{[X]} & Geprüft und negativ. \\
\end{tabular}
\end{center}

% ============================================================
\section{Einordnung: Geometrie, Algebra und was ein Körper hier ist}
% ============================================================

Die FFGFT hat zwei Schichten, die man nicht verwechseln darf.

\paragraph{Die Geometrie} ist ein Torus. Der Träger ist $T^4$ mit dem
$D_4$-Gitter (Kusszahl~24, Dok.~314), und die Physik sitzt in den
\emph{Moden} dieses Torus: den Wicklungszahlen der ebenen Wellen, die auf
ihm leben. Die eingebaute $\mathbb{Z}_3$-Symmetrie des Gitters macht daraus
den Orbifold $T^4/\mathbb{Z}_3$ mit neun Fixpunkten. Was ein Teilchen ist,
welche Masse es trägt, welcher Sektor ihm zukommt -- das sind Aussagen
über Moden und ihre Orbits unter $\mathbb{Z}_3$.

\paragraph{Die Algebra} sind endliche Körper: $GF(9)$ als Grundkörper,
$GF(27)$ mit seinen 26 Einheiten in acht Dreierorbits, beide parallel in
$GF(729)$ (Dok.~338, 341). Der Frobenius $x\mapsto x^3$ ist die
Galois-Gruppe dieser Körper, und die Sektorpaarung $k\leftrightarrow-k$ ist
seine Wirkung.

\paragraph{Wie die Schichten zusammenhängen.} Die Körper \emph{sind} nicht
die Moden. Sie sind die Eigenwertkörper, Zerlegungskörper und
Koeffizientenkörper, in denen die Modenoperatoren leben, und sie
entscheiden, welche Moden \emph{existieren können}: eine Mode ist eine
Elementordnung in einer Matrixalgebra über $GF(3^k)$, und Lagrange
erlaubt nur Ordnungen, die die Gruppenordnung $|GF(3^k)^*|$ teilen.
Deshalb gibt es Ordnung~5 in $G(3)$ ($5\mid80$), Ordnung~26 aber erst in
$G(6)$ ($13\nmid80$, Dok.~341). Kurz: \emph{Die Geometrie stellt die
Moden; die Galois-Struktur der Körper sagt, welche davon zulässig sind.}
Die Körper sind modenbedingt in genau diesem Sinn -- sie werden durch die
Moden aufgerufen, nicht umgekehrt.

\paragraph{Warum das den Vergleich ordnet.} OPH hat dieselbe Zweiteilung
mit anderen Einträgen: die Geometrie ist eine Sphäre aus Ikosaedern, die
Moden sind die Blöcke des Portraums unter $A_5$, und die Algebra ist der
Zahlkörper $\mathbb{Q}(\sqrt5)$ mit seiner Galois-Gruppe $\sqrt5\mapsto-\sqrt5$.
Auch dort entscheidet die Galois-Struktur, welche Moden über dem
Grundkörper existieren: über $\mathbb{Q}$ nur der 6D-Block, nie ein
einzelnes Triplett. Der Vergleich in diesem Dokument ist deshalb ein
Vergleich der \emph{Algebren, die die Moden klassifizieren} -- und die
beiden Algebren berühren sich exakt einmal: bei der Reduktion modulo~3,
die $\mathbb{Q}(\sqrt5)$ auf $GF(9)$ und die Galois-Konjugation auf den
Frobenius abbildet (Satz~\ref{thm:frob}). Die Geometrien dagegen
berühren sich nur in der Tetraedergruppe und trennen sich an der
Fünfzähligkeit (Abschnitt~\ref{sec:d4}).

% ============================================================
\section{Anlass und Abgrenzung}
% ============================================================

Auf der IPI-Liste wurde im September 2026 ein Trägermodell vorgestellt
(Observer Patch Holography, im Folgenden \emph{OPH}), dessen Hardware ein
endliches Netz von Zwölf-Port-Ikosaedern ist und dessen Software eine
einzige lokale Regel: an jeder der dreißig Kanten (\emph{Nähte}) werden
die beiden Port-Lesungen durch ihr arithmetisches Mittel ersetzt. Der Autor
stellte ein Reproduzierbarkeitspaket mit exakter rationaler Arithmetik
bereit (300 Updates, Seeds 7 und 11, unabhängiger Verifier).

Dieses Dokument nimmt zu den physikalischen Ansprüchen von OPH keine
Stellung. Es beantwortet drei enge, exakt prüfbare Fragen:
\begin{enumerate}[nosep]
\item Welche Struktur hat das Spektrum des Ikosaeder-Trägers, und was folgt
      aus seiner Galois-Gruppe für die Unterscheidbarkeit der beiden
      dreidimensionalen Blöcke?
\item Wie verhält sich $\mathbb{Q}(\sqrt5)$ zum FFGFT-Grundkörper $GF(9)$?
\item Welcher Teil der Ikosaedergruppe $A_5$ ist mit dem $D_4$-Gitter und
      dem Orbifold $T^4/\mathbb{Z}_3$ verträglich?
\end{enumerate}
Die Verdopplung $6=3\oplus3'$ (Dok.~285), die Einbettung $C_3<A_5$
(Dok.~293), der Negativbefund $5\nmid|{\rm Aut}(D_4)|$ (Dok.~314) und die
$G(3)/G(6)$-Schwelle (Dok.~341) sind die FFGFT-Bezugspunkte.

% ============================================================
\section{Das Spektrum des Ikosaeder-Trägers}
% ============================================================

\begin{theorem}[Spektrum, \BM]\label{thm:spektrum}
Sei $A$ die Adjazenzmatrix des Ikosaedergraphen (12 Ecken, 30 Kanten,
5-regulär) und $L=5I-A$ sein Laplace-Operator. Dann
\[
  \operatorname{spec}(A)=\{5^{(1)},\ \sqrt5^{(3)},\ -1^{(5)},\ -\sqrt5^{(3)}\},\qquad
  \operatorname{spec}(L)=\{0^{(1)},\ (5-\sqrt5)^{(3)},\ 6^{(5)},\ (5+\sqrt5)^{(3)}\}.
\]
Der einzige rationale nichttriviale Laplace-Eigenwert ist $6$ (Multiplizität~5).
Der vom Spektrum erzeugte Körper ist $\mathbb{Q}(\sqrt5)$, Grad~2.
\end{theorem}

\begin{proof}
Prüfskript \texttt{spectrum\_icosahedral\_carrier.py}, SymPy, exakt; die
Nahtliste stammt aus dem OPH-Paket, alternativ aus den kanonischen
Ecken $(0,\pm1,\pm\varphi)$ zyklisch. Das Skript spielt zusätzlich die
300 Updates von Seed~7 in rationaler Arithmetik nach und reproduziert den
im Paket protokollierten zentrierten Restwert exakt.
\end{proof}

Die Multiplizitäten $(1,3,5,3)$ sind die Dimensionen der irreduziblen
$A_5$-Darstellungen $A\oplus T_1\oplus H\oplus T_2$ (Notation der
Ikosaedergruppe; $H$ ist fünfdimensional, $T_1,T_2$ die beiden
inäquivalenten Tripletts). Auf einer Fünffachdrehung $g_5$ tragen sie die
Charakterwerte
\[
  \chi_{T_1}(g_5)=\varphi,\qquad \chi_{T_2}(g_5)=1-\varphi,\qquad
  \chi_{T_1}+\chi_{T_2}=1 .
\]

\begin{remark}[Dämpfungsraten, \BM]
Die erwartete Wirkung eines zufälligen Einzelnaht-Updates ist
$E[M]=I-L/60$ mit Blockeigenwerten $(55+\sqrt5)/60\approx0{,}954$ auf $T_1$,
$9/10$ auf $H$ und $(55-\sqrt5)/60\approx0{,}879$ auf $T_2$. $T_1$ ist der
langsamste Block; der Abstand zu $T_2$ beträgt $\sqrt5/30$
(Verhältnis $\approx1{,}085$). Das ist eine Ratenaussage, keine Symmetrieauswahl.
\end{remark}

% ============================================================
\section[Galois-Konjugation und die Verdopplung 6 = 3 + 3']{Galois-Konjugation und die Verdopplung $6=3\oplus3'$}
% ============================================================

\begin{theorem}[Galois vertauscht $T_1\leftrightarrow T_2$, \BM]\label{thm:galois}
Seien $P_{T_1},P_{T_2},P_A,P_H$ die Spektralprojektoren von $L$ (Polynome in
$L$ mit Koeffizienten in $\mathbb{Q}(\sqrt5)$). Der Galois-Automorphismus
$\sigma:\sqrt5\mapsto-\sqrt5$ von $\mathbb{Q}(\sqrt5)$ wirkt eintragsweise als
\[
  \sigma(P_{T_1})=P_{T_2},\qquad \sigma(P_{T_2})=P_{T_1},\qquad
  \sigma(P_A)=P_A,\qquad \sigma(P_H)=P_H .
\]
Der Projektor $P_{T_1}+P_{T_2}$ ist $\mathbb{Q}$-rational; $P_{T_1}$ und
$P_{T_2}$ einzeln sind es nicht. Das Minimalpolynom des $T_1$-Eigenwerts
über $\mathbb{Q}$ ist $x^2-10x+20$ mit den Wurzeln $5\mp\sqrt5$.
\end{theorem}

\begin{proof}
Prüfskript \texttt{galois\_z3\_icosahedral\_carrier.py}. Die Projektoren
werden als Lagrange-Polynome in $L$ gebildet, die Vertauschung durch
Substitution $\sqrt5\to-\sqrt5$ und Vergleich der Matrizen geprüft.
\end{proof}

\begin{corollary}[Unteilbarkeit über $\mathbb{Q}$, \BM]\label{cor:sechs}
Jede über $\mathbb{Q}$ definierte, $A_5$-invariante Struktur auf dem
Portraum -- ein Gitter, ein ganzzahliger Charakter, exakte rationale
Reparaturarithmetik -- sieht $T_1\oplus T_2$ als \emph{einen}
sechsdimensionalen Block. Die Wahl \enquote{$T_1$ ist der Raum} ist eine
Wahl des Vorzeichens von $\sqrt5$, keine Ableitung.
\end{corollary}

Das ist der Satz hinter Dok.~285: die Verdopplung $6=3\oplus3'$ ist dort
die Antwort auf die kristallographische Einschränkung (kein ikosaedrisches
Gitter in drei Dimensionen, wohl aber in sechs). Hier erscheint dieselbe
Aussage als Galois-Unteilbarkeit. Der Vorschlag von OPH, das
dreidimensionale Kontinuum sei die dichte Projektion \enquote{ganzzahliger
Rahmen aus Portlabels}, ist die Quasikristall-Projektion
$\mathbb{Z}^6\to\mathbb{R}^3$; Dok.~285 beschreibt die Ursache, OPH das Bild.

\begin{theorem}[Restriktion auf $C_3<A_5$, \BM]\label{thm:c3}
Unter einer Dreifachdrehung $g_3\in A_5$ zerfällt der Portraum in
$\mathbb{Z}_3$-Eigenräume der Dimensionen $4+4+4$; blockweise
\[
  A\to 1,\qquad T_1\to 1\oplus\omega\oplus\omega^2,\qquad
  T_2\to 1\oplus\omega\oplus\omega^2,\qquad H\to 1\oplus2\omega\oplus2\omega^2 .
\]
Der triviale Anteil hat Dimension~4, der nichttriviale Dimension~8.
\end{theorem}

Dies ist die $C_3$-in-$A_5$-Einbettung von Dok.~293 in Charakterform.
Der nichttriviale Anteil (Dimension~8) ist als abstrakter
$\mathbb{Z}_3$-Modul zwei Kopien der Orbifold-Darstellung auf $\mathbb{R}^4$;
ob darin mehr steckt als Darstellungstheorie, klärt Abschnitt~\ref{sec:d4}
negativ.

% ============================================================
\section[{Reduktion modulo 3: Z[phi]/(3) = GF(9)}]{Reduktion modulo 3: $\mathbb{Z}[\varphi]/(3)=GF(9)$}
% ============================================================

\begin{theorem}[Träge Primzahl, \EM/\BM]\label{thm:inert}
$x^2-x-1$ ist über $GF(3)$ irreduzibel ($5\equiv2$ ist kein Quadrat
modulo~3). Also ist $3$ in $\mathbb{Q}(\sqrt5)$ träge und
$\mathbb{Z}[\varphi]/(3)\cong GF(9)$.
\end{theorem}

\begin{theorem}[Galois wird Frobenius, \BM]\label{thm:frob}
Unter der Reduktion von Satz~\ref{thm:inert} geht die Galois-Konjugation
$\varphi\mapsto1-\varphi$ in den Frobenius $x\mapsto x^3$ von $GF(9)$ über:
in $GF(9)$ gilt exakt $\varphi^3=1-\varphi$.
\end{theorem}

\begin{proof}
$\varphi^2=\varphi+1$, also $\varphi^3=2\varphi+1\equiv1-\varphi\pmod3$.
Prüfskript \texttt{zphi\_mod3\_a5.py}, Abschnitte (1)--(2).
\end{proof}

Damit sind die Sektorpaarung der FFGFT (Frobenius auf $GF(27)^*$,
Dok.~343/363) und die Vertauschung $T_1\leftrightarrow T_2$ des
Ikosaeder-Trägers nicht nur analog: modulo~3 sind sie dieselbe Abbildung.

\begin{theorem}[$A_5$ über $GF(9)$, \BM]\label{thm:a5gf9}
Die sechzig Rotationsmatrizen des Ikosaeders (Einträge in
$\tfrac12\mathbb{Z}[\varphi]$) reduzieren modulo~3 auf sechzig verschiedene
Matrizen in $GL_3(GF(9))$, die ganz $M_3(GF(9))$ aufspannen: die Reduktion
ist treu und absolut irreduzibel. Eintragsweiser Frobenius der reduzierten
$T_1$-Matrizen liefert genau die reduzierten $T_2$-Matrizen.
Dagegen enthält $GL_2(GF(9))$ (Ordnung 5760, vollständig durchsucht) keine
$(2,3,5)$-erzeugte Untergruppe der Ordnung~60: $A_5$ bettet nicht in
$GL_2(GF(9))$ ein.
\end{theorem}

\begin{corollary}[Die $G(3)/G(6)$-Schwelle, \SM]
Mit $G(3)\cong M_2(GF(9))^2$ und $G(6)\cong M_8(GF(9))$ (Dok.~341) liegt
$A_5$ nicht in der Einheitengruppe von $G(3)$, aber in der von $G(6)$.
Das ist dieselbe Schwelle, an der Dok.~341 die Leptonen-Ordnung~26 findet.
Die Gründe sind formal verschieden (Brauer-Grade $\{1,3,3,4\}$ von $A_5$
in Charakteristik~3 gegen Lagrange in $GF(81)^*$); ob die Koinzidenz
Substanz hat, bleibt offen.
\end{corollary}

% ============================================================
\section[A5 gegen Aut(D4)]{$A_5$ gegen ${\rm Aut}(D_4)$}\label{sec:d4}
% ============================================================

\begin{theorem}[Reproduktion von Dok.~314, \BM]\label{thm:aut}
Die Automorphismengruppe des $D_4$-Wurzelsystems, aus den 24 Wurzeln neu
berechnet, hat Ordnung $1152$ mit Ordnungsverteilung
$(1,2,3,4,6,8,12)\mapsto(1,139,80,228,464,144,96)$ und kein Element der
Ordnung~5. Unter den 80 Elementen der Ordnung~3 sind genau 16 fixpunktfrei
auf $\mathbb{R}^4$ ($|\det(1-A)|=9$); jedes wirkt frei auf den 24 Wurzeln
(acht Orbits zu drei).
\end{theorem}

\begin{corollary}[\XM]
$A_5$ wirkt auf keinem $D_4$-Gitter: $5\nmid1152$. Die
Symmetriegruppe von OPH und die FFGFT-Gitterstruktur sind als Ganzes unverträglich.
\end{corollary}

\begin{theorem}[Der gemeinsame Kern ist $A_4$, \BM]\label{thm:a4}
Die Elemente von $A_5$, deren Matrizen im Modell $(0,\pm1,\pm\varphi)$
ganzzahlig sind (kein $\varphi$), bilden die Tetraedergruppe $A_4$
(Ordnungen $1,2,3$ mit Anzahlen $1,3,8$). $A_4$ wirkt einfach transitiv
auf den zwölf Ports -- der Portsatz ist ein $A_4$-Torsor -- und bettet als
$\operatorname{diag}(M,1)$ in $W(D_4)<{\rm Aut}(D_4)$ ein.
\end{theorem}

Das präzisiert Dok.~285 (\enquote{$C_3<A_5$}): nicht nur die Dreizähligkeit,
sondern die ganze Tetraedergruppe der Ordnung~12 -- die Portzahl -- ist
gitterverträglich; erst die Fünfzähligkeit fällt heraus.

\begin{theorem}[Ports als halbe Erstschale, \BM]
Jede Dreifachdrehung von $A_5$ wirkt frei auf den zwölf Ports (vier Orbits
zu drei). Als $\mathbb{Z}_3$-Modul ist der Portraum $4\cdot\mathbb{Z}[\mathbb{Z}_3]$,
die erste $D_4$-Schale $8\cdot\mathbb{Z}[\mathbb{Z}_3]$: die Ports sind die
eine Hälfte von \enquote{$24=12+12$} aus Dok.~314.
\end{theorem}

\begin{theorem}[Die Port-$\mathbb{Z}_3$ ist keine Orbifold-$\mathbb{Z}_3$, \XM]\label{thm:neg}
Die Dreifachdrehung, die die Ports permutiert, ist eine Drehung des
$\mathbb{R}^3$ und fixiert eine Achse; als $\operatorname{diag}(M,1)$ fixiert sie
eine Gerade in $\mathbb{R}^4$ und ist damit eines der $80-16=64$
Nicht-Orbifold-Elemente der Ordnung~3. Das ist keine Folge der Einbettung:
Für jede der 16 Orbifold-Erzeugenden und jede der 70 Wahlen von vier ihrer
acht Wurzelorbits (1120 $\mathbb{Z}_3$-stabile Zwölferhälften) wurde geprüft,
ob die zwölf Wurzeln ein Kuboktaeder ($D_3$-Hälfte) bilden oder auf einer
Inner-Product-Klasse den Ikosaedergraphen tragen. Beides tritt nie ein.
\end{theorem}

\begin{proof}
Eine $\mathbb{Z}_3$, die eine $D_3$-Hälfte stabilisiert, stabilisiert deren
Normale bis aufs Vorzeichen, bei Ordnung~3 also exakt; sie hat einen
Fixvektor und ist keine Orbifold-Erzeugende. Prüfskript
\texttt{a5\_vs\_d4.py}, Abschnitte (E)--(F).
\end{proof}

\begin{remark}[Nebenbefund, \KM]
256 der 1120 orbifold-stabilen Zwölferhälften sind Tight-Frames in
$\mathbb{R}^4$ (Gram-Eigenwert~6 mit Multiplizität~4, das 4D-Gegenstück zur
OPH-Relation $G^2=4G$ der Portvektoren). Ihre Graphenspektren enthalten
$\sqrt{13}$, nicht $\sqrt5$; der Stabilisator hat Ordnung~18 und ist nicht
transitiv. Sie tragen keine Portstruktur.
\end{remark}

% ============================================================
\section{Die Eichzuordnung von OPH}
% ============================================================

OPH ordnet dem Portmodul die Lie-Algebra $u(1)\oplus su(2)\oplus su(3)$ zu:
$u(1)$ auf $A$, $su(2)$ auf $T_1$, $su(3)$ auf $T_2\oplus H$.

\begin{theorem}[Konsistenz der Zuordnung, \EM]
$A_5$ liegt via $T_2$ in $SO(3)\subset SU(3)$. Die Adjungierte von $su(3)$
ist $3\otimes\bar3-1$; unter $SO(3)$ zerfällt sie als $3\oplus5$, unter
$A_5$ mit der $T_2$-Einbettung als $T_2\oplus H$ ($\Lambda^2T_2=T_2$,
$\operatorname{Sym}^2T_2=1\oplus H$). Ebenso ist die Adjungierte von
$su(2)$ unter der $T_1$-Einbettung $T_1$. Die Restriktion von
$u(1)\oplus su(2)\oplus su(3)$ auf $A_5$ ist also exakt
$A\oplus T_1\oplus T_2\oplus H$.
\end{theorem}

\begin{remark}[Was Axiom ist, \SM]
Die Zuordnung setzt voraus, dass die reversible Portantwort eine
zwölfdimensionale Lie-Algebra ist, auf der $A_5$ durch innere
Automorphismen wirkt, mit Zentrum genau der $A_5$-Fixgeraden. Dann ist
$11=8+3$ die einzige Zerlegung in Dimensionen kompakter einfacher Algebren.
Ohne die Zentrumsannahme hat z.\,B.\ $u(1)^3\oplus su(2)^3$ ebenfalls
Dimension~12. Korrekt ist daher \enquote{$A_5$ plus Antwortaxiom erzwingt
den Lie-Typ}, nicht \enquote{$A_5$ erzwingt den Lie-Typ}.
\end{remark}

Die FFGFT hat die Spiegellage: $su(3)$ folgt aus der $\mathbb{Z}_3$-Trialität
auf $T^4$ \BM\ (Dok.~321), $su(2)$ ist offen \SM\ (Dok.~324). Bei OPH ist
$T_1$ sauber und $8=T_2\oplus H$ der axiomatische Schritt.

% ============================================================
\section{Phasen als Einheitswurzeln: wie die Verbindung gelingt}\label{sec:phasen}
% ============================================================

Abschnitt~\ref{sec:d4} und Satz~\ref{thm:frob} scheinen sich zu
widersprechen: Die Fünfzähligkeit ist mit dem Gitter unverträglich, und
doch liegt $A_5$ nach Reduktion in $GL_3(GF(9))$. Der Widerspruch löst sich,
sobald man unterscheidet, \emph{als was} eine Symmetrie in den Träger
eintritt -- als Gittervektor-Automorphismus oder als Phase.

\begin{theorem}[Phasenkanal, \BM]\label{thm:phasen}
\begin{enumerate}[nosep]
\item $\varphi$ ist modulo~3 ein primitives Element von $GF(9)^*$:
      $\varphi^4=-1$, Ordnung~8.
\item Die $72^\circ$-Drehung hat das charakteristische Polynom
      $(x-1)(x^2-(\varphi-1)x+1)$. Modulo~3 hat der quadratische Faktor keine
      Wurzel in $GF(9)$; seine Wurzeln liegen in $GF(81)$ und haben dort
      Ordnung~5 ($5\mid80$).
\item $13\mid728=|GF(729)^*|$, aber $13\nmid80$: Ordnung~26 braucht die
      kubische Erweiterung (Dok.~341).
\item In Charakteristik~3 ist $x^3-1=(x-1)^3$: es gibt in keinem $GF(3^k)$
      eine primitive dritte Einheitswurzel. Eine $\mathbb{Z}_3$ wirkt im
      Körper unipotent, nie als Phase~$\omega$.
\end{enumerate}
\end{theorem}

\begin{proof}
Prüfskript \texttt{phasen\_einheitswurzeln.py}.
\end{proof}

\paragraph{Lesart.} Eine Phase ist eine Einheitswurzel, und Einheitswurzeln
sind genau das, was eine Körpererweiterung hinzufügt. Im Gitter ist die
Drehung um $72^\circ$ ein Automorphismus von Vektoren, und $5\nmid1152$.
Im Körper ist dieselbe Drehung nur noch ihr Eigenwert $\zeta_5$, und
$5\mid80$. Die geometrische Unverträglichkeit wird zur arithmetischen
Verträglichkeit, sobald man von Gittervektoren zu Phasen übergeht. Das ist
der Mechanismus hinter Dok.~293: $\theta=2/9$ entsteht aus $R_5$ als
Phase, nicht als Gitteroperation.

Die Erweiterungsgrade sind dabei die beiden Grundoperationen der FFGFT:
\begin{center}
\begin{tabular}{@{}llll@{}}
\toprule
Erweiterung & Grad & trägt & Gegenstück \\
\midrule
$GF(9)\to GF(81)$ & 2 & $\zeta_5$, $\varphi$, $2/\sqrt5$, Ikosaeder & Verdopplung $6=3\oplus3'$ (Dok.~285) \\
$GF(9)\to GF(729)$ & 3 & $\zeta_{13}$, Ordnung~26, Leptonen & $\mathbb{Z}_3$-Trialität (Dok.~341) \\
\bottomrule
\end{tabular}
\end{center}

Und die Dreizähligkeit selbst ist \emph{keine} Phase: In Charakteristik~3
fällt $\omega$ auf $1$ zusammen (Dok.~341: die $N$-Eigenwerte $\{-3,+3\}$
werden $\{0\}$). Deshalb ist die $\mathbb{Z}_3$ bei der FFGFT
Gitterstruktur -- Zirkulant auf den Modenorbits, Dok.~314 -- und die
Fünfzähligkeit Phase. Das ist keine Wahl, sondern erzwungen: Was in
Charakteristik~3 nicht als Einheitswurzel existieren kann, muss als
Gitteroperation kommen; was nicht ins Gitter passt, muss als
Einheitswurzel in einer Erweiterung kommen. Die beiden Kanäle sind
komplementär, und die Primzahl~3 entscheidet, welche Symmetrie in welchen
Kanal fällt. Satz~\ref{thm:neg} ist die geometrische Seite derselben
Aussage: Die Port-$\mathbb{Z}_3$ ist Gitter, die Orbifold-$\mathbb{Z}_3$
ist Gitter, aber sie sind verschiedene Gitteroperationen; die
Fünfzähligkeit, die sie unterscheidet, ist im Gitter gar nicht vorhanden.

% ============================================================
\section{Zusammenfassung}
% ============================================================

\begin{center}
\begin{tabular}{@{}lllp{4.6cm}@{}}
\toprule
 & OPH & FFGFT & Verbindung \\
\midrule
Körper & $\mathbb{Q}(\sqrt5)$ & $GF(9)$, $GF(27)\parallel GF(729)$ & $\mathbb{Z}[\varphi]/(3)=GF(9)$; Galois $\to$ Frobenius \BM \\[2pt]
Gruppe & $A_5$ & ${\rm Aut}(D_4)$, $\mathbb{Z}_3$ & $A_5\not\subset{\rm Aut}(D_4)$; Kern $A_4$ (= 12 Ports) \BM \\[2pt]
$3\oplus3'$ & $T_1$ Raum, $T_2$ Farbe & 6D-Gitter (Dok.~285) & über $\mathbb{Q}$ unteilbar \BM \\[2pt]
$\mathbb{Z}_3$-Orbits & 4 freie auf 12 & 8 freie auf 24 & Ports = halbe Erstschale \BM \\[2pt]
Port-$\mathbb{Z}_3$ & Achse in $\mathbb{R}^3$ & Orbifold fixpunktfrei & nie konjugiert \XM \\[2pt]
$G(3)/G(6)$ & $A_5\not\subset GL_2(GF(9))$, $\subset GL_3$ & Ord.~26 erst in $G(6)$ & gleiche Schwelle \SM \\
\bottomrule
\end{tabular}
\end{center}

\paragraph{Prüfskripte.} \texttt{2/python/Dok364\_Skripte/}:
\texttt{spectrum\_icosahedral\_carrier.py},
\texttt{galois\_z3\_icosahedral\_carrier.py},
\texttt{zphi\_mod3\_a5.py}, \texttt{a5\_vs\_d4.py}, \texttt{phasen\_einheitswurzeln.py}; Sammellauf
\texttt{run\_all\_364.sh}. Die Nahtliste des externen Pakets wird, falls
vorhanden, eingelesen; andernfalls wird der kanonische Ikosaeder verwendet.
Kein Registereintrag: es wird kein älteres Dokument korrigiert.
